% EXP-004: OTel Weaver Registry to wasm4pm-compat Witness Generator
% Lineage and Context Chapter

\section{Lineage and Context}
\label{sec:otel-weaver-exp-004-registry-to-witness-lineage}

\subsection{2016 Language-Model Lesson}
\label{sec:otel-weaver-exp-004-registry-to-witness-lm-lesson}

The research program begins with a 2016 experiment in local language-model text imitation.
The operative finding from that experiment was not about the quality of the imitation but
about its absence of consequence: a system that predicts the surface form of text does not
thereby inherit the causal and legal obligations that the text was used to discharge. An
imitation of a process log is not a process log. An imitation of a witness attestation is
not a witness attestation.

EXP-004 instantiates this lesson at the type-system layer. The old-regime approach to
building telemetry bridges was textual: copy field names from a convention specification
document into source code files, relying on developer vigilance to keep the two in sync.
This approach is structurally analogous to text imitation --- the code copies the shape of
the convention without being causally bound to it. When the convention changes, the code
does not change, and the failure mode is silent.

The code-generation pipeline in EXP-004 replaces textual imitation with causal derivation.
The materialized \texttt{ActivityWitness} struct is not copied from a document --- it is
forged from the same source artifact (the resolved schema JSON) that the Weaver toolchain
uses. The consequence-gap identified in the 2016 experiment is closed, at least at the
structural boundary, by making the source of truth machine-readable and the transformation
mechanically verified.

\subsection{Enterprise Process Gap}
\label{sec:otel-weaver-exp-004-registry-to-witness-enterprise-gap}

The 2016 lesson scaled to the enterprise in a specific way: distributed systems emit
telemetry from dozens of instrumented services, but the telemetry carries no
process-consequence claims. A span saying \texttt{dispatch\_payment} was emitted does not
assert that the payment was lawfully authorized. Structural validity --- all required
fields present, types correct, schema version matched --- is an engineering metric, not a
legal or operational assertion.

The doctrine file \texttt{otel-weaver/doctrine/otel-weaver-is-feedstock.md} formalizes
this gap. It draws a strict nominal boundary between schema validation (syntactic,
attribute presence, type checks, authority: systems infrastructure) and conformance
checking (temporal, activity ordering, object linkages, token replay, authority: board room
or compliance registry). The Weaver framework, however sophisticated its convention
resolution, operates entirely within the schema-validation column of this table.

EXP-004's code-generation pipeline sits at the exact boundary where the enterprise gap
manifests. The pipeline accepts structured convention metadata from the Weaver layer ---
the feedstock side of the boundary --- and forges Rust types that implement the
\texttt{Lattice} trait required by the court layer. It does not perform conformance
checking itself; it manufactures the type infrastructure that makes compile-time-verified
conformance checking possible. The gap between ``structurally valid telemetry'' and
``lawful process execution'' is not closed by EXP-004, but the plumbing that allows the
court to make that determination without hand-coded fragility is what EXP-004 provides.

\subsection{Relationship to the Chatman Equation}
\label{sec:otel-weaver-exp-004-registry-to-witness-chatman-equation}

The Chatman Equation $A = \mu(O^*)$ asserts that an autonomous agent $A$ is the measure
$\mu$ applied to the open star $O^*$ --- the totality of evidence available in a process
context. For $A$ to be a valid measure rather than an imitation, the evidence surface
$O^*$ must be structured, typed, and causally grounded. Raw telemetry is not $O^*$; it is
feedstock that must pass through a typed admission boundary before entering the evidence
surface.

EXP-004 contributes to the Chatman Equation by ensuring that the feedstock-to-evidence
projection is not a hand-written approximation. The \texttt{ActivityWitness} struct
materialized by the generator carries three fields --- \texttt{instance\_id},
\texttt{witness\_id}, and \texttt{witness\_hash} --- that correspond directly to the
attributes declared as \texttt{required} in the \texttt{process.pi.activity} group of
\texttt{process\_pi.yaml}. These are not arbitrary fields; they are the minimum typed
assertions that a process execution event must carry to be admissible as evidence in a
conformance court. The generator enforces this minimum at compile time.

The BLAKE3 hash field (\texttt{process.pi.witness.hash}) is particularly significant in
the equation. A witness hash seals a specific transition trace, making it tamper-evident.
Without this seal, the evidence surface is mutable and the measure $\mu$ is undefined.
EXP-004 ensures that every materialized witness type carries this seal as a required,
typed field, not an optional annotation.

\subsection{Relationship to Receipts and Replay}
\label{sec:otel-weaver-exp-004-registry-to-witness-receipts}

The receipt chain in this research program provides cryptographic provenance for
manufacturing steps: a receipt attests that a specific artifact was produced by a specific
process at a specific time. EXP-004 both depends on and contributes to this infrastructure.

The \texttt{generated\_witnesses.rs} artifact carries an embedded timestamp comment
(\texttt{Generated at: 2026-06-01T10:10:51-07:00}) as its sole provenance record. This is
a weak form of receipt --- it attests \emph{when} the file was materialized but not
\emph{how} (which generator version, which input schema, which operator). The absence of a
cryptographic receipt file specific to EXP-004 is a documented gap
(Section~\ref{sec:otel-weaver-exp-004-registry-to-witness-gaps}).

The replay relationship operates in the other direction: the \texttt{Lattice} trait
implementation in the materialized code is itself a replay mechanism. The
\texttt{partial\_cmp} method defines a partial order over witness states where an empty
\texttt{witness\_hash} is the bottom element (no evidence) and a non-empty 64-character
BLAKE3 hex string is the top element (sealed evidence). The \texttt{join} operation
resolves compatible witnesses by identity and resolves incompatible witnesses (distinct
non-empty hashes representing divergent execution paths) by returning the lattice top ---
a convention that signals contradiction rather than silently selecting one path. This
lattice structure is what allows \texttt{wasm4pm} to replay conformance decisions over
a partial order of witness states rather than over a flat sequence of events.
